\documentclass{article}
\usepackage[utf8]{inputenc}
\usepackage{geometry}
\usepackage{enumitem}
\usepackage{hyperref}
\usepackage{xcolor}
\usepackage{graphicx}
\usepackage{array}
\usepackage{multirow}
\usepackage{amsmath, amssymb}
\usepackage{chemfig}
\usepackage{mhchem}
\usepackage{tikz}
\usepackage{setspace}
\usepackage{soul}
\usepackage{mathtools}
\usepackage{booktabs}
\usepackage{chemformula}
\usepackage{siunitx}
\usetikzlibrary{positioning, calc}
\usepackage{longtable}
\geometry{margin=1in}
\setlength{\parskip}{0.5em}
\usepackage{cancel}

\begin{document}

\begin{center}
    {\LARGE\textbf{SI Session Planning Form}}
\end{center}

\
\noindent
\begin{flushleft}
\textbf{SI Leader:} Brandon Cobb\\
\textbf{Session Week:} 9\\
\textbf{Session\#:} 9\\
\textbf{Instructor:} Professor Tramontozzi\\
\textbf{Old Plans:}\href{https://macombedu.sharepoint.com/:b:/s/ASC-SupplementalInstructionEmbeddedTutoring/EeSYqCwT_KhMolgJvXbVEjgBOyF_DevvS8BgCm1zPkQQJg?e=QJTt4o}{SharePoint}\\
\textbf{Session Goal:} Students will have reviewed stoichiometry and balancing chemical equations, applying concepts such as formulas, atomic masses, Avogadro’s number, and mole ratios to calculate quantities of substances and relate them to reactants or products in a reaction. They will also apply gas laws and the kinetic molecular theory to determine relationships between moles, pressure, volume, temperature, and molecular speed.
\end{flushleft}

\vspace{1em}
\section*{\underline{Opener}}
{\centering\large\textbf{\hl{Take Attendance}}\par}\vspace{-0.25\baselineskip}
\noindent
\begin{tabular}{|p{4.2cm}|p{9.8cm}|}
\hline
\textbf{Collaborative Learning Techniques} & Group Discussion\\
\hline
\textbf{Learning Strategy} & Lecture Review\\
\hline
\textbf{Time} & 12:00---12:10\\
\hline
\textbf{Materials} & Notes, homework, whiteboard, markers, ready to speak, periodic table \\
\hline
\textbf{Lecture/Textbook/Old Plans/Slide References} & Class roster, Chapter 6: Chemical calculations: Formula masses, moles and chemical equations. Chapter 7: Gases, Liquids and Solids \\
\hline
\end{tabular}


\vspace{0.5cm}
{\large\textbf{Definitions}}
\\
\begin{tabular}{|p{0.9\linewidth}|}
\hline
\begin{minipage}[t]{0.9\linewidth}
\begin{enumerate}
    \item\textbf{Group Discussion:} One person in the group is asked to present on a topic or review material for the group and then lead the discussion for the group. This person should not be the regular group leader.
   \item\textbf{Lecture Review:} As a group summarize the lecture from the previous class. You may have to provide prompts for the students. For example, "The first concept discussed was Civil Liberties adn Public Policy, what did the profressor highlight regarding this?" You may want to ask them to try summarizing without looking at their notes; however, if they are having a difficult time remembering, tell them to refer to their notes.
\end{enumerate}
\end{minipage} \\
\hline
\end{tabular}
\newpage
\noindent
{\large\textbf{Instructions (please provide a\hl{DETAILED} activity breakdown below (and/or attached}}\\

\begin{tabular}{|p{0.9\linewidth}|}
\hline
\vspace{0.2cm}
\textbf{Lecture Review:}
\begin{enumerate}
\item During te first 10-15 minutes of the SI session. have the students summarize the most recent lecture, or have them identify the key words from that lecture.
\item Give students three minutes to find support in their lecture notes for a given generalization.
\item Have the students predict the direction of future lectures based upon the past lectures.
\item Have the student arrange terms fgrom lecture and text into a structured outline.
\item Reinforce new terms or important information by using clearly constructed handouts (can be complete or nearly complete at the beginning of the term but should gradually require more and more filling in as the group becomes more accustomed to working together.
\item Review material from previous sessions and lectures.
\item Take a couple of minutes at the end of the SI session to summarize the main idea convered during a session. Asl the students to help summarize.
\item Have students write a one paragraph summary of the lecture. List the new vocabulary terms introduced with this lecture.
\item Formulate potential exam questions based on the main ideas from the lecture.
\item Formulate pontential answers from details in the lecture notes.
\end{enumerate} \\
\hline
\end{tabular}

\newpage
\noindent
{\large\textbf{Checking for Understanding (questions \& answers)} \\
\vspace{0.2cm}
\begin{tabular}{|p{0.9\linewidth}|}
\hline
\begin{itemize}
\item How does knowing the number of moles help predict the amount of product formed in a reaction?
\item Why is the ratio of reactants to products important in chemical reactions?
\item How do coefficients in a chemical equation relate to mole ratios?
\item How can stoichiometry be used to determine how much of one substance is needed to react with another?
\item Why is it important to consider limiting reactants in a reaction?
\item How does Avogadro’s number allow us to relate moles to the number of particles?
\item How does molar mass connect the mass of a substance to the number of moles?
\item Why do reactions with different compounds require careful attention to their specific formula masses?
\item How can stoichiometry help compare quantities of different chemicals in a reaction?
\item How does gas pressure relate to the number of gas particles and their collisions with container walls?
\item How does temperature affect the speed of gas particles?
\item Why do lighter gas molecules move faster than heavier ones at the same temperature?
\item How does increasing the volume of a gas at constant temperature affect its pressure?
\end{itemize} \\
\hline
\end{tabular}

\newpage
\noindent
{\large\textbf{Checking for Understanding (questions \& answers) continued...} \\
\vspace{0.2cm}
\begin{tabular}{|p{0.9\linewidth}|}
\hline
\begin{itemize}
\item How does increasing the number of gas particles in a fixed container affect pressure?
\item Why do real gases deviate from ideal behavior at high pressures or low temperatures?
\item How does the kinetic molecular theory explain gas diffusion?
\item How do intermolecular forces influence boiling points and melting points?
\item Why do polar molecules tend to have higher boiling points than nonpolar molecules of similar size?
\item How does hydrogen bonding affect the viscosity and surface tension of a liquid?
\item How does the strength of intermolecular forces influence evaporation and vapor pressure?
\item Why are gas particle collisions considered elastic?
\item How does the kinetic energy of gas particles relate to the temperature of the gas?
\item Why does diffusion occur faster for lighter gases compared to heavier gases?
\item How can understanding mole ratios and gas laws together help predict reaction outcomes involving gases?
\item How does the spacing between gas particles affect how closely a real gas behaves like an ideal gas?
\item How do molecular mass and temperature together affect the average speed of gas molecules?
\end{itemize} \\
\hline
\end{tabular}


\section*{\underline{Activity 1}}
\vspace{0.2cm}
\begin{tabular}{|p{4.2cm}|p{9.8cm}|}
\hline
\textbf{Collaborative Learning Techniques} & Jigsaw\\
\hline
\textbf{Learning Strategy} & Predict Test Questions \\
\hline
\textbf{Time} & 12:10---12:30\\
\hline
\textbf{Materials} & Pencil, eraser, calculator, periodic table \\
\hline
\textbf{Lecture/Textbook/Old Plans/Slide References} & Chapter 6: Chemical calculations: Formula masses, moles and chemical equations. Chapter 7: Gases, Liquids and Solids \\
\hline
\end{tabular}
\vspace{0.5cm}

{\large\textbf{Definition}}\\
\noindent
\begin{tabular}{|p{0.9\linewidth}|}
\hline
\begin{enumerate}
    \item\textbf{Jigsaw:} Jigsaws, when used properly, make the group as a whole dependent upon all of them in subgroups. Each group provides a peice of the puzzle. Group members are broken into smaller groups. Each small group works on some aspect of the same problem, question, or issue. They then share their part of the puzzle with the large group.
    \item\textbf{Predict Test Questions:}  Put students in groups of 2 or 3 and assign them to write a test question for a specific topic, ensuring that all topics have been convered. Ask students to write their questions on the board or on an overhead for discussion (would the professor ask the question? What is the answer? etc.). Students will have the benefit of learning to think like the teacher and they'll be able to see additional questions that other students have writen.
 
\end{enumerate}\\
\hline
\end{tabular}
\vspace{0.2cm}
\newpage
{\large\textbf{Instructions (please provide a\hl{DETAILED} activity breakdown below (and/or attached}}\\
\noindent
\begin{tabular}{|p{0.9\linewidth}|}
\hline
\vspace{0.2cm}
\begin{enumerate}
    \item \textbf{Distribute Worksheets:} Each student receives a worksheet containing chemistry questions. The problems cover a range of difficulty so that students can decide which ones best match their needs.
    
    \item \textbf{Student Choice of Starting Point:} Students look over the worksheet and select their own starting problem rather than being assigned a fixed order. This allows them to challenge themselves where they feel it is most beneficial.
    
    \item \textbf{Staggered Timing:} Students work in pairs. Each partner begins at a different point (for example, one may start at Problem 1 while the other starts at Problem 3). They continue by alternating every three questions. This staggered structure ensures that students will be asking for help or requesting solutions at different times, keeping the class flow balanced.
    
    \item \textbf{Independent Work:} Students attempt their chosen problems individually, writing down both reasoning and answers. Emphasis is placed on documenting thought process as well as the result.
    
    \item \textbf{Help and Solution Reveal:} After the set time for each problem expires, the SI leader either walks through the solution on the whiteboard or directs students to check the answer key. Students are encouraged to ask clarifying questions at the timed checkpoints rather than all at once.
    
    \item \textbf{Partner Collaboration:} Once a solution is reviewed, partners share their reasoning for their respective problems. They ask each other questions, offer corrections, and compare how their approaches aligned with the solution.
    
    \item \textbf{Rotation of Problems:} After every three questions, students switch to new problems, continuing the staggered order so that they stay offset from each other. This ensures that each student engages with a variety of problems across the worksheet.
    
    \item \textbf{Class-Wide Debrief:} The SI leader pauses periodically to highlight frequently asked questions, difficult concepts, or common errors observed during the staggered work. Whole-class solutions are discussed when many students identify the same challenge.
    
    \item \textbf{Reflection:} At the close, students note which problems gave them the most insight, what strategies helped, and what topics they still want clarified. This reflection helps focus future review sessions.
\end{enumerate}\\
\hline
\end{tabular}
\newpage

\noindent
{\large\textbf{Content (fully worked out problems \& answers) continued...}}\\[0.2em]
\vspace{0.2cm}

\noindent\begin{tabular}{|p{0.9\linewidth}|}
\hline
\vspace{0.2cm}
\textbf{(1)} How many moles are in 25.0~g of NaCl?\\[0.2cm]

\begin{minipage}[t]{0.85\linewidth}
\begin{enumerate}[label=\alph*.]
\item 0.427
\item 2.50
\item 0.125
\item 1.00
\end{enumerate}
\end{minipage}\\[0.2cm]
\hline
\end{tabular}

\vspace{0.2cm}
\noindent\begin{tabular}{|p{0.9\linewidth}|}
\hline
\textbf{(2)} $2\text{H}_2 + \text{O}_2 \rightarrow 2\text{H}_2\text{O}$. If 3.0~mol of H$_2$ react with excess O$_2$, how many moles of H$_2$O are formed?\\[0.2cm]

\begin{minipage}[t]{0.85\linewidth}
\begin{enumerate}[label=\alph*.]
\item 1.5
\item 3.0
\item 6.0
\item 2.0
\end{enumerate}
\end{minipage}\\[0.2cm]
\hline
\end{tabular}

\vspace{0.2cm}
\noindent\begin{tabular}{|p{0.9\linewidth}|}
\hline
\textbf{(3)} $\text{N}_2 + 3\text{H}_2 \rightarrow 2\text{NH}_3$. If 4.0~mol H$_2$ react with 2.0~mol N$_2$, which is the limiting reagent?\\[0.2cm]

\begin{minipage}[t]{0.85\linewidth}
\begin{enumerate}[label=\alph*.]
\item N$_2$
\item H$_2$
\item Neither
\item Both
\end{enumerate}
\end{minipage}\\[0.2cm]
\hline
\end{tabular}

\vspace{0.2cm}
\noindent\begin{tabular}{|p{0.9\linewidth}|}
\hline
\textbf{(4)} In the reaction $\text{C} + \text{O}_2 \rightarrow \text{CO}_2$, how many grams of CO$_2$ form when 10.0~g of carbon reacts completely?\\[0.2cm]

\begin{minipage}[t]{0.85\linewidth}
\begin{enumerate}[label=\alph*.]
\item 22.0
\item 36.7
\item 10.0
\item 44.0
\end{enumerate}
\end{minipage}\\[0.2cm]
\hline
\end{tabular}
\newpage

\noindent\begin{tabular}{|p{0.9\linewidth}|}
\hline
\textbf{(5)} $\text{Ca} + 2\text{H}_2\text{O} \rightarrow \text{Ca(OH)}_2 + \text{H}_2$. How many liters of H$_2$ are produced at 298~K and 1.50~atm from 2.00~mol Ca?\\[0.2cm]

\begin{minipage}[t]{0.85\linewidth}
\begin{enumerate}[label=\alph*.]
\item 36.5
\item 44.8
\item 24.7
\item 32.6
\end{enumerate}
\end{minipage}\\[0.2cm]
\hline
\end{tabular}

\vspace{0.2cm}
\noindent\begin{tabular}{|p{0.9\linewidth}|}
\hline
\textbf{(6)} $2\text{K} + \text{Cl}_2 \rightarrow 2\text{KCl}$. How many moles of KCl are formed from 0.500~mol Cl$_2$?\\[0.2cm]

\begin{minipage}[t]{0.85\linewidth}
\begin{enumerate}[label=\alph*.]
\item 0.250
\item 1.00
\item 0.500
\item 2.00
\end{enumerate}
\end{minipage}\\[0.2cm]
\hline
\end{tabular}

\vspace{0.2cm}
\noindent\begin{tabular}{|p{0.9\linewidth}|}
\hline
\textbf{(7)} How many molecules are present in 0.300~mol of CO$_2$?\\[0.2cm]

\begin{minipage}[t]{0.85\linewidth}
\begin{enumerate}[label=\alph*.]
\item $1.81\times10^{23}$
\item $3.01\times10^{23}$
\item $6.02\times10^{23}$
\item $9.03\times10^{23}$
\end{enumerate}
\end{minipage}\\[0.2cm]
\hline
\end{tabular}

\vspace{0.2cm}
\noindent\begin{tabular}{|p{0.9\linewidth}|}
\hline
\textbf{(8)} If 50.0~g of Fe reacts with excess O$_2$ to form Fe$_2$O$_3$, and the actual yield is 68.0~g, what is the percent yield if the theoretical yield was 80.0~g?\\[0.2cm]

\begin{minipage}[t]{0.85\linewidth}
\begin{enumerate}[label=\alph*.]
\item 85.0\%
\item 90.0\%
\item 68.0\%
\item 80.0\%
\end{enumerate}
\end{minipage}\\[0.2cm]
\hline
\end{tabular}

\newpage

\noindent\begin{tabular}{|p{0.9\linewidth}|}
\hline
\textbf{(9)} $\text{Zn} + 2\text{HCl} \rightarrow \text{ZnCl}_2 + \text{H}_2$. How many grams of ZnCl$_2$ form when 0.500~mol of Zn reacts completely?\\[0.2cm]

\begin{minipage}[t]{0.85\linewidth}
\begin{enumerate}[label=\alph*.]
\item 34.0
\item 68.0
\item 136
\item 17.0
\end{enumerate}
\end{minipage}\\[0.2cm]
\hline
\end{tabular}

\vspace{0.2cm}
\noindent\begin{tabular}{|p{0.9\linewidth}|}
\hline
\textbf{(10)} How many oxygen atoms are in 0.200~mol of Al$_2$O$_3$?\\[0.2cm]

\begin{minipage}[t]{0.85\linewidth}
\begin{enumerate}[label=\alph*.]
\item $3.61\times10^{23}$
\item $1.20\times10^{23}$
\item $6.02\times10^{23}$
\item $7.23\times10^{23}$
\end{enumerate}
\end{minipage}\\[0.2cm]
\hline
\end{tabular}

\vspace{0.2cm}
\noindent\begin{tabular}{|p{0.9\linewidth}|}
\hline
\textbf{(11)} $2\text{H}_2 + \text{O}_2 \rightarrow 2\text{H}_2\text{O}$. If 4.0~mol H$_2$ react with 2.0~mol O$_2$, which is limiting?\\[0.2cm]

\begin{minipage}[t]{0.85\linewidth}
\begin{enumerate}[label=\alph*.]
\item H$_2$
\item O$_2$
\item Neither
\item Both
\end{enumerate}
\end{minipage}\\[0.2cm]
\hline
\end{tabular}

\vspace{0.2cm}
\noindent\begin{tabular}{|p{0.9\linewidth}|}
\hline
\textbf{(12)} If the theoretical yield of NH$_3$ from $\text{N}_2 + 3\text{H}_2 \rightarrow 2\text{NH}_3$ is 10.0~g, and the actual yield is 7.80~g, what is the percent yield?\\[0.2cm]

\begin{minipage}[t]{0.85\linewidth}
\begin{enumerate}[label=\alph*.]
\item 78.0\%
\item 100\%
\item 82.0\%
\item 70.0\%
\end{enumerate}
\end{minipage}\\[0.2cm]
\hline
\end{tabular}

\newpage

\noindent\begin{tabular}{|p{0.9\linewidth}|}
\hline
\textbf{(13)} How many moles of oxygen gas are needed to completely burn 2.0~mol CH$_4$? ($\text{CH}_4 + 2\text{O}_2 \rightarrow \text{CO}_2 + 2\text{H}_2\text{O}$)\\[0.2cm]

\begin{minipage}[t]{0.85\linewidth}
\begin{enumerate}[label=\alph*.]
\item 1.0
\item 2.0
\item 4.0
\item 8.0
\end{enumerate}
\end{minipage}\\[0.2cm]
\hline
\end{tabular}

\vspace{0.2cm}
\noindent\begin{tabular}{|p{0.9\linewidth}|}
\hline
\textbf{(14)} $2\text{Na} + \text{Cl}_2 \rightarrow 2\text{NaCl}$. If 5.00~mol Na react with 1.50~mol Cl$_2$, what is the limiting reagent?\\[0.2cm]

\begin{minipage}[t]{0.85\linewidth}
\begin{enumerate}[label=\alph*.]
\item Na
\item Cl$_2$
\item Neither
\item Both
\end{enumerate}
\end{minipage}\\[0.2cm]
\hline
\end{tabular}

\vspace{0.2cm}
\noindent\begin{tabular}{|p{0.9\linewidth}|}
\hline
\textbf{(15)} If 0.750.~mol CaCO$_3$ decomposes to $\text{CaO} + \text{CO}_2$, how many moles of CO$_2$ are produced?\\[0.2cm]

\begin{minipage}[t]{0.85\linewidth}
\begin{enumerate}[label=\alph*.]
\item 1.50
\item 0.750
\item 0.375
\item 0.250
\end{enumerate}
\end{minipage}\\[0.2cm]
\hline
\end{tabular}

\vspace{0.2cm}
\noindent\begin{tabular}{|p{0.9\linewidth}|}
\hline
\textbf{(16)} $\text{Zn} + \text{S} \rightarrow \text{ZnS}$. If 3.00~mol Zn react with 2.00~mol S, which is the limiting reagent?\\[0.2cm]

\begin{minipage}[t]{0.85\linewidth}
\begin{enumerate}[label=\alph*.]
\item Zn
\item S
\item Neither
\item Both
\end{enumerate}
\end{minipage}\\[0.2cm]
\hline
\end{tabular}

\newpage

\noindent\begin{tabular}{|p{0.9\linewidth}|}
\hline
\textbf{(17)} If 5.00~mol of O$_2$ react with excess H$_2$, how many grams of H$_2$O are formed? ($2\text{H}_2 + \text{O}_2 \rightarrow 2\text{H}_2\text{O}$)\\[0.2cm]

\begin{minipage}[t]{0.85\linewidth}
\begin{enumerate}[label=\alph*.]
\item 90.0
\item 45.0
\item 180
\item 22.4
\end{enumerate}
\end{minipage}\\[0.2cm]
\hline
\end{tabular}

\vspace{0.2cm}
\noindent\begin{tabular}{|p{0.9\linewidth}|}
\hline
\textbf{(18)} $2\text{Al} + 3\text{Cl}_2 \rightarrow 2\text{AlCl}_3$. If 2.0~mol Al react with 4.0~mol Cl$_2$, what is the limiting reagent?\\[0.2cm]

\begin{minipage}[t]{0.85\linewidth}
\begin{enumerate}[label=\alph*.]
\item Al
\item Cl$_2$
\item Neither
\item Both
\end{enumerate}
\end{minipage}\\[0.2cm]
\hline
\end{tabular}

\vspace{0.2cm}
\noindent\begin{tabular}{|p{0.9\linewidth}|}
\hline
\textbf{(19)} $2\text{H}_2\text{O}_2 \rightarrow 2\text{H}_2\text{O} + \text{O}_2$. If 6.00~mol H$_2$O$_2$ decompose, how many moles of O$_2$ form?\\[0.2cm]

\begin{minipage}[t]{0.85\linewidth}
\begin{enumerate}[label=\alph*.]
\item 3.00
\item 6.00
\item 12.0
\item 1.50
\end{enumerate}
\end{minipage}\\[0.2cm]
\hline
\end{tabular}

\vspace{0.2cm}
\noindent\begin{tabular}{|p{0.9\linewidth}|}
\hline
\textbf{(20)} What is the percent yield if 10.0~g of product are expected but only 8.20~g are obtained?\\[0.2cm]

\begin{minipage}[t]{0.85\linewidth}
\begin{enumerate}[label=\alph*.]
\item 82.0\%
\item 92.0\%
\item 80.0\%
\item 88.0\%
\end{enumerate}
\end{minipage}\\[0.2cm]
\hline
\end{tabular}

\newpage

\noindent\begin{tabular}{|p{0.9\linewidth}|}
\hline
\textbf{(21)} $2\text{HCl} + \text{Mg} \rightarrow \text{MgCl}_2 + \text{H}_2$. How many moles of H$_2$ are produced from 5.00~mol of HCl?\\[0.2cm]

\begin{minipage}[t]{0.85\linewidth}
\begin{enumerate}[label=\alph*.]
\item 5.00
\item 2.50
\item 10.0
\item 1.00
\end{enumerate}
\end{minipage}\\[0.2cm]
\hline
\end{tabular}

\vspace{0.2cm}
\noindent\begin{tabular}{|p{0.9\linewidth}|}
\hline
\textbf{(22)} Complete the combustion reaction of methane and calculate how many grams of CO$_2$ are produced from 0.500~mol CH$_4$: \\[0.2cm]
\hspace{1em} $\text{CH}_4 + \text{O}_2 \rightarrow \text{CO}_2 + \text{H}_2\text{O}$ \\[0.2cm]

\begin{minipage}[t]{0.85\linewidth}
\begin{enumerate}[label=\alph*.]
\item 11.0
\item 22.0
\item 44.0
\item 16.0
\end{enumerate}
\end{minipage}\\[0.2cm]
\hline
\end{tabular}

\vspace{0.2cm}
\noindent\begin{tabular}{|p{0.9\linewidth}|}
\hline
\textbf{(23)} $2\text{Na} + 2\text{H}_2\text{O} \rightarrow 2\text{NaOH} + \text{H}_2$. If 4.00~mol Na react, how many moles of H$_2$ are released?\\[0.2cm]

\begin{minipage}[t]{0.85\linewidth}
\begin{enumerate}[label=\alph*.]
\item 1.00
\item 2.00
\item 4.00
\item 8.00
\end{enumerate}
\end{minipage}\\[0.2cm]
\hline
\end{tabular}

\vspace{0.2cm}
\noindent\begin{tabular}{|p{0.9\linewidth}|}
\hline
\textbf{(24)} In the reaction $\text{CaCO}_3 \rightarrow \text{CaO} + \text{CO}_2$, if 50.0~g of CaCO$_3$ produce 25.0~g of CaO, what is the percent yield if 28.0~g was expected?\\[0.2cm]

\begin{minipage}[t]{0.85\linewidth}
\begin{enumerate}[label=\alph*.]
\item 89.3\%
\item 75.0\%
\item 92.0\%
\item 100\%
\end{enumerate}
\end{minipage}\\[0.2cm]
\hline
\end{tabular}

\newpage

\noindent\begin{tabular}{|p{0.9\linewidth}|}
\hline
\textbf{(25)} $\text{Fe}_2\text{O}_3 + 3\text{CO} \rightarrow 2\text{Fe} + 3\text{CO}_2$. If 2.00~mol Fe$_2$O$_3$ react with 6.00~mol CO, what mass of Fe is formed assuming 100\% yield?\\[0.2cm]

\begin{minipage}[t]{0.85\linewidth}
\begin{enumerate}[label=\alph*.]
\item 56.0
\item 112
\item 224
\item 168
\end{enumerate}
\end{minipage}\\[0.2cm]
\hline
\end{tabular}

\vspace{0.2cm}
\noindent\begin{tabular}{|p{0.9\linewidth}|}
\hline
\textbf{(26)} A 425~mL cylinder contains nitrogen at $2.50\times10^{2}$~kPa and 22.0$^\circ$C. What pressure (in atm) will result if the gas is moved to a 1.00~L container and heated to 125$^\circ$C, assuming moles constant?\\[0.2cm]

\begin{minipage}[t]{0.85\linewidth}
\begin{enumerate}[label=\alph*.]
\item 0.564~atm
\item 1.41~atm
\item 0.282~atm
\item 2.24~atm
\end{enumerate}
\end{minipage}\\[0.2cm]
\hline
\end{tabular}

\vspace{0.2cm}
\noindent\begin{tabular}{|p{0.9\linewidth}|}
\hline
\textbf{(27)} 0.375~mol of an ideal gas occupies $6.50\times10^{2}$~mL at 27.0$^\circ$C. Using the ideal gas law, what is the pressure in mmHg? \\[0.2cm]

\begin{minipage}[t]{0.85\linewidth}
\begin{enumerate}[label=\alph*.]
\item 575~mmHg
\item 760.~mmHg
\item 430.~mmHg
\item $1.08\times10^{4}$~mmHg
\end{enumerate}
\end{minipage}\\[0.2cm]
\hline
\end{tabular}

\vspace{0.2cm}
\noindent\begin{tabular}{|p{0.9\linewidth}|}
\hline
\textbf{(28)} A sealed syringe has 8.00~mL of gas at 295~K and 750.~torr. If the plunger is pulled to 24.0~mL and temperature falls to 270.~K, what is the new pressure (in kPa)?\\[0.2cm]

\begin{minipage}[t]{0.85\linewidth}
\begin{enumerate}[label=\alph*.]
\item 125~kPa
\item 188~kPa
\item 62.5~kPa
\item 30.5~kPa
\end{enumerate}
\end{minipage}\\[0.2cm]
\hline
\end{tabular}

\newpage

\noindent\begin{tabular}{|p{0.9\linewidth}|}
\hline
\textbf{(29)} 1.50~g of helium is confined in a 2.00~L flask at 350.~K. Using the ideal gas law, what is the pressure in atm? \\[0.2cm]

\begin{minipage}[t]{0.85\linewidth}
\begin{enumerate}[label=\alph*.]
\item 0.650.~atm
\item 5.39~atm
\item 0.325~atm
\item 1.45~atm
\end{enumerate}
\end{minipage}\\[0.2cm]
\hline
\end{tabular}

\vspace{0.2cm}
\noindent\begin{tabular}{|p{0.9\linewidth}|}
\hline
\textbf{(30)} A gas at 1.20~atm and 300~K fills a 5.00~L vessel. If the temperature is increased to 450.~K and pressure is allowed to equilibrate at constant moles, what is the final volume (L)?\\[0.2cm]

\begin{minipage}[t]{0.85\linewidth}
\begin{enumerate}[label=\alph*.]
\item 7.50~L
\item 3.33~L
\item 5.00~L
\item 11.3~L
\end{enumerate}
\end{minipage}\\[0.2cm]
\hline
\end{tabular}

\vspace{0.2cm}
\noindent\begin{tabular}{|p{0.9\linewidth}|}
\hline
\textbf{(31)} A sample of gas has a volume of 2.50~L at a pressure of 1.20~atm. If the pressure is increased to 3.00~atm at constant temperature, what is the new volume of the gas? (Boyle’s Law — indirect relationship)\\[0.2cm]

\begin{minipage}[t]{0.85\linewidth}
\begin{enumerate}[label=\alph*.]
\item 1.00~L
\item 2.10.~L
\item 0.950~L
\item 3.75~L
\end{enumerate}
\end{minipage}\\[0.2cm]
\hline
\end{tabular}

\vspace{0.2cm}
\noindent\begin{tabular}{|p{0.9\linewidth}|}
\hline
\textbf{(32)} A gas mixture in a 10.0~L tank exerts total pressure 4.00~atm. The partial pressures are $P_A=1.50$~atm and $P_B=0.75$~atm. What is $P_C$ and how many moles of C are present at 25.0$^\circ$C?\\[0.2cm]

\begin{minipage}[t]{0.85\linewidth}
\begin{enumerate}[label=\alph*.]
\item $P_C=1.75$~atm, $n_C=7.22\times10^{-1}$~mol
\item $P_C=1.75$~atm, $n_C=1.63$~mol
\item $P_C=1.25$~atm, $n_C=0.722$~mol
\item $P_C=1.75$~atm, $n_C=3.61\times10^{-1}$~mol
\end{enumerate}
\end{minipage}\\[0.2cm]
\hline
\end{tabular}

\newpage

\noindent\begin{tabular}{|p{0.9\linewidth}|}
\hline
\textbf{(33)} A gas initially at 640.~mmHg and 20.0$^\circ$C is compressed to 300~mL from 750.~mL while temperature changes to 60.0$^\circ$C. What is final pressure in mmHg?\\[0.2cm]

\begin{minipage}[t]{0.85\linewidth}
\begin{enumerate}[label=\alph*.]
\item $1.12\times10^{3}$~mmHg
\item 512~mmHg
\item 290.~mmHg
\item $1.50\times10^{3}$~mmHg
\end{enumerate}
\end{minipage}\\[0.2cm]
\hline
\end{tabular}

\vspace{0.2cm}
\noindent\begin{tabular}{|p{0.9\linewidth}|}
\hline
\textbf{(34)} 0.250~mol of oxygen gas is collected at 740.~mmHg and 18.0$^\circ$C in a 2.00~L flask. What would be the pressure at 0.500~mol in same flask and same T? (assume ideal)\\[0.2cm]

\begin{minipage}[t]{0.85\linewidth}
\begin{enumerate}[label=\alph*.]
\item 1480.~mmHg
\item 740.~mmHg
\item 370.~mmHg
\item 1110.~mmHg
\end{enumerate}
\end{minipage}\\[0.2cm]
\hline
\end{tabular}

\vspace{0.2cm}
\noindent\begin{tabular}{|p{0.9\linewidth}|}
\hline
\textbf{(35)} A gas at 1.00~atm and 298~K is allowed to expand isothermally from 2.00~L to 6.00~L. What is final pressure in atm?\\[0.2cm]

\begin{minipage}[t]{0.85\linewidth}
\begin{enumerate}[label=\alph*.]
\item 3.00~atm
\item 0.333~atm
\item 0.667~atm
\item 1.00~atm
\end{enumerate}
\end{minipage}\\[0.2cm]
\hline
\end{tabular}

\vspace{0.2cm}
\noindent\begin{tabular}{|p{0.9\linewidth}|}
\hline
\textbf{(36)} A sample of gas weighing 4.80~g exerts 0.900~atm in a 2.50~L container at 300.~K. What is the molar mass of the gas? \\[0.2cm]

\begin{minipage}[t]{0.85\linewidth}
\begin{enumerate}[label=\alph*.]
\item 58.2~g~mol$^{-1}$
\item 120.~g~mol$^{-1}$
\item 52.5~g~mol$^{-1}$
\item 16.0~g~mol$^{-1}$
\end{enumerate}
\end{minipage}\\[0.2cm]
\hline
\end{tabular}

\newpage

\noindent\begin{tabular}{|p{0.9\linewidth}|}
\hline
\textbf{(37)} At 1.00~atm and 298~K, a gas sample has density 0.860~g/L. What is its molar mass?\\[0.2cm]

\begin{minipage}[t]{0.85\linewidth}
\begin{enumerate}[label=\alph*.]
\item 19.4~g~mol$^{-1}$
\item 24.7~g~mol$^{-1}$
\item 29.0~g~mol$^{-1}$
\item 256~g~mol$^{-1}$
\end{enumerate}
\end{minipage}\\[0.2cm]
\hline
\end{tabular}

\vspace{0.2cm}
\noindent\begin{tabular}{|p{0.9\linewidth}|}
\hline
\textbf{(38)} A 150.~mL sample of gas at 25$^\circ$C and 755~mmHg contains $2.00\times10^{21}$ molecules. What is Avogadro's number implied by these measurements? (students must convert units and compute $N_A$)\\[0.2cm]

\begin{minipage}[t]{0.85\linewidth}
\begin{enumerate}[label=\alph*.]
\item $6.02\times10^{23}$
\item $5.12\times10^{23}$
\item $7.30\times10^{23}$
\item $3.01\times10^{23}$
\end{enumerate}
\end{minipage}\\[0.2cm]
\hline
\end{tabular}

\vspace{0.2cm}
\noindent\begin{tabular}{|p{0.9\linewidth}|}
\hline
\textbf{(39)} A gas occupies 500.~mL at 25$^\circ$C. What volume will it occupy at 100.$^\circ$C if the pressure remains constant? \\[0.2cm]

\begin{minipage}[t]{0.85\linewidth}
\begin{enumerate}[label=\alph*.]
\item 583~mL
\item 625~mL
\item 665~mL
\item 740.~mL
\end{enumerate}
\end{minipage}\\[0.2cm]
\hline
\end{tabular}

\vspace{0.2cm}
\noindent\begin{tabular}{|p{0.9\linewidth}|}
\hline
\textbf{(40)} A gas initially at 1.20~atm, 298~K and 250.~mL is heated to 350.~K and allowed to expand until pressure is 1.00~atm. What is the final volume (mL)?\\[0.2cm]

\begin{minipage}[t]{0.85\linewidth}
\begin{enumerate}[label=\alph*.]
\item 330.~mL
\item 250.~mL
\item 290.~mL
\item 400~mL
\end{enumerate}
\end{minipage}\\[0.2cm]
\hline
\end{tabular}

\newpage

\noindent\begin{tabular}{|p{0.9\linewidth}|}
\hline
\textbf{(41)} A gas is stored in a container at 2.00~atm and 20.0$^\circ$C. If the temperature is increased to 80.0$^\circ$C at constant volume, what will be the new pressure? \\[0.2cm]

\begin{minipage}[t]{0.85\linewidth}
\begin{enumerate}[label=\alph*.]
\item 2.35~atm
\item 2.60.~atm
\item 2.80.~atm
\item 3.00~atm
\end{enumerate}
\end{minipage}\\[0.2cm]
\hline
\end{tabular}

\vspace{0.2cm}
\noindent\begin{tabular}{|p{0.9\linewidth}|}
\hline
\textbf{(42)} A cylinder contains 0.850~mol of gas at 2.50~atm. What is the temperature (K) if the volume is 15.0~L?\\[0.2cm]

\begin{minipage}[t]{0.85\linewidth}
\begin{enumerate}[label=\alph*.]
\item 210.~K
\item 298~K
\item 425~K
\item 512~K
\end{enumerate}
\end{minipage}\\[0.2cm]
\hline
\end{tabular}

\vspace{0.2cm}
\noindent\begin{tabular}{|p{0.9\linewidth}|}
\hline
\textbf{(43)} A sealed vessel at 1.00~atm and 300.~K contains $3.00\times10^{22}$ molecules. If vessel is evacuated and recharged with the same number of molecules at 600.~K, what will be the new pressure (atm) in the same vessel?\\[0.2cm]

\begin{minipage}[t]{0.85\linewidth}
\begin{enumerate}[label=\alph*.]
\item 0.500~atm
\item 1.00~atm
\item 2.00~atm
\item 4.00~atm
\end{enumerate}
\end{minipage}\\[0.2cm]
\hline
\end{tabular}

\vspace{0.2cm}
\noindent\begin{tabular}{|p{0.9\linewidth}|}
\hline
\textbf{(44)} True or False: Gas particles are in constant, random motion.\\[0.2cm]

\begin{minipage}[t]{0.85\linewidth}
\begin{enumerate}[label=\alph*.]
\item True
\item False
\end{enumerate}
\end{minipage}\\[0.2cm]
\hline
\end{tabular}

\newpage

\noindent\begin{tabular}{|p{0.9\linewidth}|}
\hline
\textbf{(45)} True or False: Gas pressure is caused by collisions of particles with the container walls.\\[0.2cm]

\begin{minipage}[t]{0.85\linewidth}
\begin{enumerate}[label=\alph*.]
\item True
\item False
\end{enumerate}
\end{minipage}\\[0.2cm]
\hline
\end{tabular}

\vspace{0.2cm}
\noindent\begin{tabular}{|p{0.9\linewidth}|}
\hline
\textbf{(46)} Increasing the temperature of a gas increases the average:\\[0.2cm]

\begin{minipage}[t]{0.85\linewidth}
\begin{enumerate}[label=\alph*.]
\item Pressure only
\item Molecular speed
\item Volume only
\item Mass
\end{enumerate}
\end{minipage}\\[0.2cm]
\hline
\end{tabular}

\vspace{0.2cm}
\noindent\begin{tabular}{|p{0.9\linewidth}|}
\hline
\textbf{(47)} True or False: Lighter gas particles move faster than heavier ones at the same temperature.\\[0.2cm]

\begin{minipage}[t]{0.85\linewidth}
\begin{enumerate}[label=\alph*.]
\item True
\item False
\end{enumerate}
\end{minipage}\\[0.2cm]
\hline
\end{tabular}

\vspace{0.2cm}
\noindent\begin{tabular}{|p{0.9\linewidth}|}
\hline
\textbf{(48)} Gas particles collide with each other:\\[0.2cm]

\begin{minipage}[t]{0.85\linewidth}
\begin{enumerate}[label=\alph*.]
\item Inelastically, losing energy each time
\item Perfectly elastically, no net energy loss
\item Only with container walls
\item Only at high pressure
\end{enumerate}
\end{minipage}\\[0.2cm]
\hline
\end{tabular}

\newpage

\noindent\begin{tabular}{|p{0.9\linewidth}|}
\hline
\textbf{(49)} True or False: Real gases behave more like ideal gases at low pressure.\\[0.2cm]

\begin{minipage}[t]{0.85\linewidth}
\begin{enumerate}[label=\alph*.]
\item True
\item False
\end{enumerate}
\end{minipage}\\[0.2cm]
\hline
\end{tabular}

\vspace{0.2cm}
\noindent\begin{tabular}{|p{0.9\linewidth}|}
\hline
\textbf{(50)} A gas sample occupies 1.20~L when it contains 0.050~mol of gas. What will its volume be if 0.200~mol of gas is added at constant temperature and pressure? \\[0.2cm]

\begin{minipage}[t]{0.85\linewidth}
\begin{enumerate}[label=\alph*.]
\item 2.40.~L
\item 3.60.~L
\item 4.80.~L
\item 6.00~L
\end{enumerate}
\end{minipage}\\[0.2cm]
\hline
\end{tabular}

\vspace{0.2cm}
\noindent\begin{tabular}{|p{0.9\linewidth}|}
\hline
\textbf{(51)} True or False: Gas particles are so small that their volume is negligible compared to the volume of the container.\\[0.2cm]

\begin{minipage}[t]{0.85\linewidth}
\begin{enumerate}[label=\alph*.]
\item True
\item False
\end{enumerate}
\end{minipage}\\[0.2cm]
\hline
\end{tabular}

\vspace{0.2cm}
\noindent\begin{tabular}{|p{0.9\linewidth}|}
\hline
\textbf{(52)} True or False: All gases at the same temperature have the same average kinetic energy.\\[0.2cm]

\begin{minipage}[t]{0.85\linewidth}
\begin{enumerate}[label=\alph*.]
\item True
\item False
\end{enumerate}
\end{minipage}\\[0.2cm]
\hline
\end{tabular}

\newpage

\noindent\begin{tabular}{|p{0.9\linewidth}|}
\hline
\textbf{(53)} Increasing the number of gas particles in a container at constant volume and temperature:\\[0.2cm]

\begin{minipage}[t]{0.85\linewidth}
\begin{enumerate}[label=\alph*.]
\item Increases pressure
\item Decreases molecular speed
\item Lowers temperature
\item Decreases collisions
\end{enumerate}
\end{minipage}\\[0.2cm]
\hline
\end{tabular}

\vspace{0.2cm}
\noindent\begin{tabular}{|p{0.9\linewidth}|}
\hline
\textbf{(54)} True or False: Increasing volume at constant temperature decreases gas pressure.\\[0.2cm]

\begin{minipage}[t]{0.85\linewidth}
\begin{enumerate}[label=\alph*.]
\item True
\item False
\end{enumerate}
\end{minipage}\\[0.2cm]
\hline
\end{tabular}

\vspace{0.2cm}
\noindent\begin{tabular}{|p{0.9\linewidth}|}
\hline
\textbf{(55)} True or False: Gas particles attract each other significantly in an ideal gas.\\[0.2cm]

\begin{minipage}[t]{0.85\linewidth}
\begin{enumerate}[label=\alph*.]
\item True
\item False
\end{enumerate}
\end{minipage}\\[0.2cm]
\hline
\end{tabular}

\vspace{0.2cm}
\noindent\begin{tabular}{|p{0.9\linewidth}|}
\hline
\textbf{(56)} True or False: The faster a gas particle moves, the more kinetic energy it has.\\[0.2cm]

\begin{minipage}[t]{0.85\linewidth}
\begin{enumerate}[label=\alph*.]
\item True
\item False
\end{enumerate}
\end{minipage}\\[0.2cm]
\hline
\end{tabular}

\newpage

\noindent\begin{tabular}{|p{0.9\linewidth}|}
\hline
\textbf{(57)} Which of the following increases when temperature increases for a gas at constant volume?\\[0.2cm]

\begin{minipage}[t]{0.85\linewidth}
\begin{enumerate}[label=\alph*.]
\item Pressure
\item Molecular mass
\item Intermolecular forces
\item Volume
\end{enumerate}
\end{minipage}\\[0.2cm]
\hline
\end{tabular}

\vspace{0.2cm}
\noindent\begin{tabular}{|p{0.9\linewidth}|}
\hline
\textbf{(58)} A gas has a volume of 3.00~L at 760.~mmHg and 300.~K. What will be its volume at 600.~mmHg and 350.~K? \\[0.2cm]

\begin{minipage}[t]{0.85\linewidth}
\begin{enumerate}[label=\alph*.]
\item 4.50.~L
\item 5.00~L
\item 5.25~L
\item 6.00~L
\end{enumerate}
\end{minipage}\\[0.2cm]
\hline
\end{tabular}

\vspace{0.2cm}
\noindent\begin{tabular}{|p{0.9\linewidth}|}
\hline
\textbf{(59)} True or False: Gas particle collisions are perfectly elastic, meaning no kinetic energy is lost.\\[0.2cm]

\begin{minipage}[t]{0.85\linewidth}
\begin{enumerate}[label=\alph*.]
\item True
\item False
\end{enumerate}
\end{minipage}\\[0.2cm]
\hline
\end{tabular}

\vspace{0.2cm}
\noindent\begin{tabular}{|p{0.9\linewidth}|}
\hline
\textbf{(60)} True or False: Gas particles are very far apart compared to their size.\\[0.2cm]

\begin{minipage}[t]{0.85\linewidth}
\begin{enumerate}[label=\alph*.]
\item True
\item False
\end{enumerate}
\end{minipage}\\[0.2cm]
\hline
\end{tabular}

\newpage
{\large\textbf{Checking for Understanding (general conceptual questions)}}\\
\noindent \begin{tabular}{|p{0.9\linewidth}|}
\hline
\vspace{0.2cm}

\begin{enumerate}
    \item \textbf{Targeted focus through self-selection:} When students choose which questions to tackle from the worksheet, they naturally gravitate toward areas they feel less confident about, ensuring the group addresses gaps in understanding.
    \item \textbf{Peer modeling of problem-solving:} Seeing a peer work through a problem on the board exposes students to different ways of setting up and explaining a solution, clarifying concepts for observers.
    \item \textbf{Reflection on understanding:} Students must decide whether they can solve a question themselves or need help, prompting metacognition and highlighting what they truly understand versus what requires clarification.
    \item \textbf{Immediate correction of misconceptions:} Peers correcting or expanding an explanation in real time prevents misunderstandings from becoming ingrained.
    \item \textbf{Balancing independence and collaboration:} Students first work independently on a question, then collaborate during discussion, ensuring individual accountability while leveraging group problem-solving.
    \item \textbf{Distinguishing conceptual vs. procedural questions:} Students can classify questions by identifying key ideas versus stepwise calculations, determining whether a question emphasizes big-picture reasoning or step-by-step execution.
    \item \textbf{Connecting big concepts to details:} Rotating through different student-selected questions helps the group link overarching concepts to multiple specific examples, reinforcing understanding of the framework behind the problems.
    \item \textbf{Verbalizing reasoning:} Explaining steps out loud forces students to articulate assumptions, logic, and chemical principles, improving clarity and communication of ideas.
    \item \textbf{Learning from multiple approaches:} Different students may solve the same problem uniquely, exposing peers to alternative methods or perspectives and broadening their problem-solving toolkit.
    \item \textbf{Application of theory to practice:} Working collaboratively on worksheet problems allows students to connect abstract chemical principles to concrete calculations or predictions.
    \item \textbf{Instructor insight:} Questions students select or ask reveal where common difficulties lie, helping instructors or SI leaders provide targeted guidance.
    \item \textbf{Strengthening long-term understanding:} Actively solving, explaining, and correcting questions in real time consolidates memory and understanding more effectively than passive review.
\end{enumerate}
\hline
\end{tabular}

\newpage
\noindent
\newpage
\noindent
\section*{\underline{Activity 2}}
\begin{tabular}{|p{4.2cm}|p{9.8cm}|}
\hline
\textbf{Collaborative Learning Techniques} & Clusters \\
\hline
\textbf{Learning Strategy} & Escape Room \\
\hline
\textbf{Time} & 12:35---12:50 \\
\hline
\textbf{Materials} & A periodic table, the textbook, readiness to work as a team or share ideas \\
\hline
\textbf{Lecture/Textbook/Old Plans/Slide References} & Chapter 6: Chemical calculations: Formula masses, moles and chemical equations. Chapter 7: Gases, Liquids and Solids \\
\hline
\end{tabular}

\vspace{0.5cm}
\noindent
\begin{tabular}{|p{0.9\linewidth}|}
\hline
{\large\textbf{Definitions}}
\begin{enumerate}
    \item \textbf{Escape Room:} Students enter a themed, time-bound challenge where they must solve a series of puzzles, riddles, or tasks to “escape” the room. Each puzzle is designed to test understanding of specific course concepts, requiring critical thinking, teamwork, and creativity. Students can expect to encounter hidden clues, coded messages, or problem-solving stations that build on one another; progress usually depends on collaboration and discussion. The SI leader will circulate to provide subtle hints, clarify misconceptions, and ensure students stay on track without giving away answers. This activity promotes peer-to-peer learning, reinforces content in an engaging way, and encourages students to communicate their reasoning while practicing time management under pressure.
\end{enumerate} \\
\hline
\end{tabular}
\newpage

\vspace{0.5cm}
\noindent
{\large \textbf{Instructions (detailed activity breakdown below)}} \\

\begin{tabular}{|p{0.9\linewidth}|}
\hline
\begin{enumerate}
    \item \textbf{Objective:} Students will collaboratively solve chemistry problems covering \textbf{stoichiometry, gas laws, KMT, and intermolecular forces}. Each correct response earns points toward an overall team score. The goal is to promote discussion, problem-solving, and application of multiple concepts in a dynamic setting.

    \item \textbf{Group Size:} Pairs (2 students per team)

    \item \textbf{Setup (Teacher Preparation):}
    \begin{itemize}
        \item Arrange \textbf{5–6 stations}, each focused on one major topic area:
        \begin{itemize}
            \item \textbf{Station 1 – Stoichiometry}
            \item \textbf{Station 2 – Gas Laws}
            \item \textbf{Station 3 – Dimensional analysis}
            \item \textbf{Station 4 – Review}
            \item \textbf{(Optional) Station 5 – Hint / Checkpoint Station}
        \end{itemize}
        \item At each station, place a set of laminated question cards (multiple choice and true/false). Each card includes:
        \begin{itemize}
            \item A unique question number and topic label.
            \item The possible answer choices (A–E or T/F).
            \item A mapping line showing which station each answer leads to (to simulate “pathways” through the escape room).
        \end{itemize}
        \item Provide each pair of students with a \textbf{team slip} (half-sheet) labeled with their team name and space to record up to 3 verified answers.
    \end{itemize}
\end{enumerate}
\\
\hline
\end{tabular}
\newpage

\vspace{0.5cm}
\noindent
{\large \textbf{Instructions (detailed activity breakdown below) conitnued...}} \\

\begin{tabular}{|p{0.9\linewidth}|}
\hline
\begin{enumerate}
    \item \textbf{Facilitation:}
    \begin{itemize}
        \item Pairs may begin at \textbf{any station of their choice}.
        \item Each pair selects one question card, discusses the problem, and writes their chosen answer \textbf{on the back of the card}.
        \item Students then show their work and reasoning to the instructor or SI leader for quick verification.
        \item If the answer is correct, they record that problem on their team slip (problem number and answer) and may proceed to any other station of their choice.
        \item Pairs are encouraged to move freely between stations, aiming to correctly solve at least \textbf{three problems} during the activity period.
        \item After completing at least three problems, teams submit their slips to the instructor for scoring.
    \end{itemize}

    \item \textbf{Scoring:}
    \begin{itemize}
        \item Each correctly solved problem = \textbf{1 point}.
        \item Teams with the highest total score win or “escape” first.
        \item Bonus points may be awarded for clear reasoning or teamwork.
    \end{itemize}

    \item \textbf{Rules:}
    \begin{itemize}
        \item Teams must work collaboratively—no switching partners mid-activity.
        \item All work must be shown before a response is accepted.
        \item Only verified answers written on the team slip count toward the final score.
        \item Communication between pairs is encouraged, but direct sharing of answers is not permitted.
    \end{itemize}

    \item \textbf{Wrap-Up:}
    \begin{itemize}
        \item After all slips are collected, review several challenging questions as a class.
        \item Discuss common reasoning errors and strategies used by top-scoring pairs.
        \item Reinforce conceptual links between stations (e.g., gas behavior and intermolecular forces).
    \end{itemize}
\end{enumerate}
\\
\hline
\end{tabular}



\newpage
\noindent
{\large\textbf{Checking for Understanding (questions \& answers)}}\\[0.2em]
\noindent
\begin{tabular}{|p{0.9\linewidth}|}
\hline
\textbf{Questions:} \\
\begin{enumerate}
    \item \textbf{Targeted clarification:} Allowing students to bring questions from the lecture lets them focus on concepts they found most confusing.
    \item \textbf{Peer modeling of explanations:} Observing how classmates summarize or interpret lecture content helps students see alternative ways to understand the material.
    \item \textbf{Reflection on comprehension:} Discussing lecture points collaboratively encourages students to distinguish what they grasp from what needs further clarification.
    \item \textbf{Immediate correction of misconceptions:} Peers can correct or expand on explanations in real time, preventing misunderstandings from solidifying.
    \item \textbf{Balancing independent thinking and discussion:} Students first reflect individually, then contribute to group discussion, promoting both accountability and collaborative learning.
    \item \textbf{Identifying conceptual vs. procedural points:} Students can highlight which parts of the lecture are big-picture ideas versus detailed steps or examples.
    \item \textbf{Connecting concepts across topics:} Discussing multiple lecture points together helps students link overarching principles to specific details.
    \item \textbf{Verbalizing reasoning:} Explaining ideas out loud strengthens students’ ability to communicate and reason about the material clearly.
    \item \textbf{Learning from diverse perspectives:} Hearing multiple interpretations or summaries of the same lecture point exposes students to alternative approaches.
    \item \textbf{Application of theory:} Collaborative discussion encourages students to connect lecture content to examples, problems, or real-world applications.
    \item \textbf{Instructor insight:} Observing which lecture points students question most helps the instructor identify common difficulties and adjust instruction.
    \item \textbf{Reinforcing long-term understanding:} Actively reviewing and discussing lecture material together solidifies comprehension better than passive note review.
\end{enumerate} \\
\hline  
\end{tabular}

\newpage
\noindent
\section*{\underline{Closer}}
\begin{tabular}{|p{4.2cm}|p{9.8cm}|}
\hline
\textbf{Collaborative Learning Techniques} & Group Survey\\
\hline
\textbf{Learning Strategy} & Brain dump\\
\hline
\textbf{Time} & 12:50---12:55\\
\hline
\textbf{Materials} & Index cards, writing utensils\\
\hline
\textbf{Lecture/Textbook/Old Plans/Slide References} & Chapters 1-7 \\
\hline
\end{tabular}

\vspace{0.5cm}
\noindent
\begin{tabular}{|p{0.9\linewidth}|}
\hline
{\large\textbf{Definitions}}
\begin{enumerate}
   \item\textbf{Group survey:} Each group member is surveyed to discover their position on an issue, problem or topic. This process ensures that each member of the group is allowed to offer or state their point of view.
   \item\textbf{Brain dump:} This is a great closer for students to sum up their knowledge on the topic(s) from the session. Have the students take out a blank piece of paper /Google Doc and write/type("dump") everything they know relevant from the discussion that day. The SI Leader can look at their work to assess the students' understanding of the course material and have them compare it to each other. It is important to at least discuss with the students so the SI Leader knows the students' comfort levels with the material. 
\end{enumerate}
\\ \hline
\end{tabular}

\newpage
\noindent
{\large\textbf{Instructions (please provide a\hl{DETAILED} activity breakdown below (and/or attached}}

\begin{tabular}{|p{0.9\linewidth}|}
\hline
\begin{enumerate}[itemsep=0.3em]
   \item Students are put into small groups and are asked to write down their confidence on a single notecard without names.
   \item The cards are evidence of the group's confidence in the content.
\end{enumerate}
\\ \hline
\end{tabular}

\vspace{0.5cm}
\noindent
{\large\textbf{Content (fully worked out problems \& answers)}

\begin{tabular}{|p{0.9\linewidth}|}
\hline
\begin{itemize}
   \item No new worked problems; brain dump.
\end{itemize}
\\ \hline
\end{tabular}

\vspace{0.5cm}
\noindent
{\large\textbf{Checking for Understanding (questions \& answers) continued...}

\begin{tabular}{|p{0.9\linewidth}|}
\hline
\begin{itemize}
   \item Review confidence levels
   \item Interpret the variety of confidence and their distribution in groups.
   \item Get feedback for what they like and don't like
   \item Check in if the time is suitable for the students
\end{itemize}
\\ \hline
\end{tabular}
\newpage
\section*{Plan Checklist}

\begin{tabular}{|p{0.9\linewidth}|}
\hline
\begin{itemize}
\item[$\Box$] Variety of Collaborative Learning Techniques and Learning Strategies
\item[$\Box$] Chapter/slide\# where information is coming from
\item[$\Box$] Included timestamps (and buffer time as needed)
\item[$\Box$] Details on how leader will break up students
\item[$\Box$] Checking for Understanding questions
\item[$\Box$] Fully worked out problems and examples for each activity
\item[$\Box$] Necessary links required for online implementation (if applicable)
\item[$\Box$] Source material correctly cited (if applicable)
\item[$\Box$] Plan is uploaded to Teams
\end{itemize}
\\ \hline
\end{tabular}
\newpage

\end{document}